\begin{abstract}
Affective computing is rapidly transitioning from controlled laboratory settings to in-the-wild applications, driven by the ubiquity of webcams and advances in computer vision. However, existing web-based physiological sensing methods predominantly assume a static, context-independent mapping between physiological arousal and affective states (e.g., ``high arousal equals stress''). This assumption fundamentally ignores the ``Invariance Breaking'' phenomenon, where identical physiological patterns can support divergent subjective states (e.g., Threat vs. Challenge) depending on cognitive appraisal and environmental context. In this paper, we propose a privacy-preserving, fully client-side web sensing framework (Camerala) designed to capture this context-dependency through Single-Case Experimental Design (SCED). We validate the system with a rigorous 10-session longitudinal study ($N=1$, 30 blocks, 600 trials) across two distinct ecological environments: a controlled University Laboratory and a naturalistic Home setting. 

Our results reveal a significant interaction between environmental context and psychophysiological reactivity. A 2-way ANOVA confirmed that while ``Threat'' stressors induced significant behavioral performance enhancement (Reaction Time acceleration) in the University setting ($F(1,396) = 9.30, p = 0.002$), this effect was significantly dampened and effectively neutralized in the Home environment ($Interaction: F(1,396) = 4.47, p = 0.035$). Bayesian analysis provided strong evidence for the threat-challenge distinction in the University context (BF$_{10}=66.7$) but inconclusive evidence in the Home context (BF$_{10}=1.7$). These findings suggest that the ``Safe Haven'' of the home environment acts as a powerful moderator of psychophysiological reactivity, buffering the impact of cognitive stressors. Investigating these intra-individual dynamics requires longitudinal, high-frequency sampling that is impossible with cross-sectional group designs. We discuss the implications of these findings for the design of next-generation, context-aware affective computing systems that move beyond context-agnostic arousal detection.
\end{abstract}

\begin{IEEEkeywords}
Web-based Sensing, Single-Case Experimental Design, Invariance Breaking, rPPG, Affective Computing, Ecological Validity, Context Dependency, Polyvagal Theory.
\end{IEEEkeywords}
